%% ma: L11-B17-0059
%% lop: 11
%% chuong: 5
%% bai: 17
%% dang: 11.17.3
%% loai: TN4
%% muc: VD
%% dap_an: "C"
%% nguon: "tai-lieu-hoc-tap-mon-toan-lop-11-theo-chuong-trinh-giao-duc-pho-thong-2018.pdf (Lớp Toán Thầy Duy, Phạm Lê Duy, THPT Chu Văn An), Chương 3 Bài 3, Phần 1, Câu 39 (tr.411)"
%% kien_thuc: "Lớp 11 Bài 17 (hàm số liên tục)"
%% trang_thai: cho-duyet
%% ly_do: "Dạng mới đề xuất (11.17.x), chờ giáo viên duyệt dạng; tài liệu gốc không có lời giải mẫu, đáp án do kho tự giải và kiểm bằng sympy."
%% ghi_chu: "Bảng đáp án gốc ghi B, nhưng B đúng (phương trình có đúng 3 nghiệm); mệnh đề sai là C (chỉ có 1 nghiệm lớn hơn 2). Đáp án do kho tự giải; cần giáo viên đối chiếu."
\begin{ex}
Cho phương trình $2x^3-8x-1=0$ $(1)$. Khẳng định nào sau đây \textbf{sai}?
\choice{Phương trình không có nghiệm lớn hơn $3$.}{Phương trình có đúng $3$ nghiệm phân biệt.}{\True Phương trình có $2$ nghiệm lớn hơn $2$.}{Phương trình có nghiệm trong khoảng $(-5;-1)$.}
\loigiai{\sloppy Đặt $f(x)=2x^3-8x-1$: $f(-2)=-1<0$, $f(-1)=5>0$, $f(0)=-1<0$, $f(2)=-1<0$, $f(3)=29>0$. Do đó có nghiệm trong $(-2;-1)$, $(-1;0)$ và $(2;3)$: đúng $3$ nghiệm phân biệt (B đúng, D đúng vì $(-2;-1)\subset(-5;-1)$). Chỉ có $1$ nghiệm lớn hơn $2$ (nằm trong $(2;3)$), nên C sai; và vì $f(x)>0$ với $x\ge3$ nên không có nghiệm lớn hơn $3$ (A đúng).}
\end{ex}
